\scanpage{46}
% L1–2, 2)
2) $\underline A\xrightarrow{T}\underline A/\underline C=\underline B$

Conditions équivalentes sur $T$ i.e. sur $\underline C$

% L3–4
\textcircled{a} $T$ admet un foncteur adjoint $S$ :
\[
\operatorname{Hom}(Tx,y)\overset{\sim}{\unsure{\mathord{-}|\to}}\operatorname{Hom}(x,Sy)
\]

% L5–9
\textcircled{b} $\forall\, x\in\underline A$, il existe un plus grand sous-objet $x'$ de $x$ qui est $\in\underline C$, \uline{et} si ce dernier est nul, il existe un monomorphisme de $x$ dans un objet $\underline C$-fermé

% L10–12
\textcircled{c} $\forall\,x\in\mathrm{Ob}\,\underline A$, il existe un hom
\[
x\longrightarrow\tilde x
\]
qui est un $\underline C$-isom, avec $\tilde x$ $\underline C$-fermé.

% L13–14, addition between the two lists
\inserted{Alors le hom en c) est unique à isom unique près, c'est aussi $x\longrightarrow STx=Lx$}

% L15–16, 1)
1) $\underline A\overset{T}{-}\underline A/\underline C=\underline B$, $x\in\mathrm{Ob}\,\underline A$. Conditions équivalentes sur $x$ (pour $\underline C$ i.e. $T$ \unsure{donné|donnés|donnée})

% L17
\textcircled{a} $\forall\,x'\in A,\ \operatorname{Hom}(x',x)\xrightarrow{\sim}\operatorname{Hom}(Tx',Tx)$

% L18–19
\textcircled{b} $\forall$ $x'-x''$ un $\underline C$-isom, on a
\[
\operatorname{Hom}(x'',x)\xrightarrow{\sim}\operatorname{Hom}(x',x)
\]

% L18–20, lower left
% slanted
\margin{\unsure{i.e.} pour tt $\underline C$-isom \unsure{$u$} \illegible{1} \unsure{$x''$}, est un isom \uline{surjectif}, \uline{\unsure{non}}, \unsure{et|i.e.} \deleted{\unsure{est}} \unsure{$\exists$} un inverse à droite}

% L20–21, torn lower edge
\textcircled{c} Tout sous objet de $x$ qui est $\in\underline C$ est nul, et \inserted{pour} tt \unsure{$y\supset x$|$y>x$|$x\subset x'$} \unsure{tel que} $\unsure{y/x|x'/x}\in\underline C$ \deleted{\unsure{est} \illegible{1}} $x$ est facteur direct !

% archive sheet number: 45
